\documentclass[12pt]{article}
\usepackage[utf8]{inputenc}
\usepackage[T1]{fontenc}
\usepackage{hyperref}
\title{Ozolotepec Zapotec Language Resources for K-12 Curriculum Development}
\author{Compiled by Claude Sonnet 5.5 (claude-sonnet-5-5) for Andr\'as Kornai}
\date{}
\begin{document}
\maketitle

\section{Introduction}
Ozolotepec Zapotec (ISO 639-3: zao; Glottocode ozol1235) is an Otomanguean language of southern Oaxaca, Mexico. Wikipedia gives about 6,500 speakers, citing the 1990 census. It is partially intelligible with Cuixtla and Loxicha Zapotec. Little is available, and no school materials were found.

\section{Primary Educational Resources}
No primers, textbooks or school curricula were found in the searches.

\section{Linguistic and Reference Materials}
Wikipedia article, which lists dialects and related varieties: \url{https://en.wikipedia.org/wiki/Ozolotepec_Zapotec}. No dictionary or grammar was found in the searches.

\section{Religious and Community Resources}
Search listings point to a New Testament in Ozolotepec Zapotec through Faith Comes By Hearing and Scripture Earth. I did not see those pages themselves, so no direct links are given.

\section{Digital and Audio Resources}
Global Recordings Network offers audio-visual Bible lessons (40 sections) with MP3 downloads and MP4 slideshows: \url{https://www.globalrecordings.net/en/language/zao}

\section{Curriculum Development Context}
The language is likely still used by the community, but materials are very scarce. A curriculum would need new documentation and orthography work with speakers. Check which town's variety is being targeted, since several towns named Ozolotepec speak different Zapotec languages.

\end{document}